\label{sec:introduction}

Open-seat elections are, by every empirical measure, more competitive than races
featuring a sitting incumbent. \citet{cox1996legislative} document that incumbency
advantage---the additional vote share a sitting member earns relative to an open-seat
candidate of the same party---averaged 6--9 percentage points in the 1980s and 1990s.
\citet{jacobson2006politics} shows that open-seat races attract higher-quality
challengers, defined as candidates with prior elected office experience, who mount
more effective campaigns and reduce the typical advantage of the majority party.
\citet{cain1987incumbency} find that open-seat races in marginal districts are the
primary mechanism by which partisan tides translate into seat changes in the House
of Representatives.

Recognising these dynamics, incumbent legislators treat open-seat creation as one of
the primary threats of redistricting. A mapmaker who draws a district that contains
no incumbent's home---whether by inadvertence or by design---creates an open seat in
which the incumbent's party has no built-in advantage. An algorithmic mapmaker who
has no information about incumbent locations will create open seats whenever the
geometric optimisation places a district boundary in a way that excludes all incumbents'
home tracts from a particular district. The rate at which this occurs is the
incumbency-neutral baseline for open-seat creation.

This paper derives that baseline, applies it to NC, WI, and TX, and compares the
algorithmic open-seat count to enacted maps and to available neutral comparators
(commission-drawn and court-drawn maps). We also develop a \textit{retirement incentive
model} that assesses the pressures each incumbent faces under an algorithmic plan:
incumbents who are paired with stronger opponents or whose district partisan lean
shifts substantially face higher retirement probability, and the aggregate retirement
pressure is higher under algorithmic than enacted plans.

\subsection{Contribution}

This paper makes three contributions. First, we derive the analytical expected
open-seat count under random redistricting, accounting for geographic concentration
of incumbent residences. Second, we test whether \textsc{bisect} algorithmic maps
and enacted maps fall within this baseline. Third, we construct a simple retirement
incentive model that translates open-seat counts, pairing outcomes, and partisan
composition changes into estimates of the retirement pressure each individual incumbent
faces.

\subsection{States and Data}

We study NC ($k=14$), WI ($k=8$), and TX ($k=38$) using 2020 Census tract data,
FEC 2022 candidate filing addresses geocoded to census tracts, and VEST 2020
presidential precinct returns. All \textsc{bisect} results use seed\,$=42$ and
are marked with $^\dagger$.
